\begin{example}[有限区间的 Robin 束缚态阈值]
\label{ex:phys-ghs-finite-robin}
半直线的束缚态条件是 $\gamma<0$。现在把粒子限制在 $(0,L)$，保留 $u'(0)=\gamma u(0)$，另一端取 $u(L)=0$。定义域为
\[
\Dom H_\gamma=\{u\in H^2(0,L):u'(0)=\gamma u(0),\ u(L)=0\},
\quad H_\gamma=-\frac{\hbar^2}{2\mu}u''.
\]
求负能束缚态存在的充要条件、能量方程与归一化系数，并与半直线极限比较。

\textbf{解。} Green 边界型在两端消失，伴随定义域的任意端点迹要求同样两条条件，故这一实现自伴。

负能量写作 $E=-\hbar^2\kappa^2/(2\mu)$，$\kappa>0$。满足右端条件的解只能取 $u=C\sinh[\kappa(L-x)]$，左端给出
\begin{equation}
\gamma=-\kappa\coth(\kappa L).
\label{eq:phys-ghs-finite-robin-root}
\end{equation}
令 $t=\kappa L$。函数 $t\coth t$ 从 $1$ 严格增加到无穷，因为
\[
\frac{\dd}{\dd t}(t\coth t)
=\frac{\sinh t\cosh t-t}{\sinh^2t}>0;
\]
分子在零点为零，导数为 $2\sinh^2t>0$。所以恰有一个负能态的充要条件是
\begin{equation}
\gamma<-\frac1L.
\label{eq:phys-ghs-finite-robin-threshold}
\end{equation}
其归一化系数由
\[
|C|^{-2}=\int_0^L\sinh^2[\kappa(L-x)]\dd x
=\frac{\sinh(2\kappa L)}{4\kappa}-\frac L2
\]
确定。阈值 $\gamma=-1/L$ 时，归一化零能态为 $u_0=\sqrt{3/L^3}(L-x)$。

也可从二次型看见这个阈值：
\[
\frac{2\mu}{\hbar^2}\langle u,H_\gamma u\rangle
=\int_0^L|u'|^2\dd x+\gamma|u(0)|^2.
\]
因 $u(L)=0$，Cauchy--Schwarz 不等式给 $|u(0)|^2\le L\int|u'|^2$，故 $-1/L\le\gamma<0$ 时二次型仍非负；$\gamma\ge0$ 时亦然。正能态满足 $\gamma=-k\cot(kL)$，与负能方程在零能阈值连续相接。

取固定负 $\gamma$ 并令 $L\to\infty$，\cref{eq:phys-ghs-finite-robin-root} 给出 $\kappa\to-\gamma$，恢复半直线的能量。有限区间的另一端提高了束缚态形成所需的边界吸引强度，因而不能直接照搬半直线的判据。
\end{example}
